
These cross-section measurements are presented in ``reconstructed'' quantities.  This avoids the ill-posed problem of trying to correct for detector smearing to present a result in true kinematics (known as unfolding~\cite{DAgostini:1994fjx}).
Instead, theoretical predictions must be smeared to compare to the data; this process called forward-folding and has been used by  T2K~\cite{Koch2019}.

To simplify the procedure for others wishing to compare predictions with our data, a single smearing matrix is provided for each measurement.  Resulting systematic uncertainties in the resolution and efficiency are estimated and included with the final data. This means that the data are corrected for an effective efficiency, as a function of a reconstructed variable.

The smearing matrix is calculated as

\begin{equation}
S_{ij} = N_{ij}^{\text{sel}} / N_{j}^{\text{sel}},
\end{equation}

%\noindent
%effectively meaning each column of the smearing matrix is normalized to 1, so that the total number of events is conserved when smearing.
\noindent
where $N_{ij}^{\text{sel}}$ is the number of selected events in reconstructed bin $i$, which come from true bin $j$, and $ N_{j}^{\text{sel}}$ is the total number of selected events from true bin $j$.
Defined in this way, a binned histogram multiplied by this smearing matrix leads to a smeared distribution with the same normalization.

In order for the data to be compared to such a prediction, they must be corrected for inefficiency.  We define an effective efficiency, $\tilde{\epsilon}_i$ which can be applied as a function of reconstructed variables:

\begin{equation}
\label{eq:eff_smear}
\tilde{\epsilon}_i = \frac{ \sum_{j=1}^{M} S_{ij}N^\text{sel}_j}{ \sum_{j=1}^{M} S_{ij}N^\text{gen}_j}, \quad 
\end{equation}

\noindent
where $N^\text{sel}_j$ is defined above, $N^\text{gen}_j$ is the number of generated signal events in true bin $j$, and $M$ is the total number of true bins used.

It is known that there are limitations to this method.  The primary limitation here is that uncertainties on the resolution are not fully encapsulated, but the uncertainty on the resolution in this particular measurement is negligible.

To determine the final differential cross sections, both signal and background event rates, as well as efficiencies, are binned as a function of each variable; the single differential cross section in bin $i$ is then calculated as (using $p_\mu$ as an example):

\begin{equation}
\label{eq:xsec_differential}
\begin{split}
\left ( \frac{d\sigma}{dp_\mu} \right )_i &= \frac{N_i - B_i}{\tilde{\epsilon}_i \cdot N_\text{target} \cdot \Phi_{\nu_\mu} \cdot (\Delta p_\mu)_i}, \\
%\left ( \frac{d\sigma}{d\cos\theta_\mu} \right )_i &= \frac{N_i - B_i}{\epsilon_i \cdot N_{target} \cdot \Phi_{\nu_\mu} \cdot (\Delta \cos\theta_\mu)_i}
\end{split}
\end{equation}

\noindent
where $N_i$, $B_i$, and $\tilde{\epsilon_i}$ are the the number of candidate events (Sec.~\ref{sec:event_sel}), the expected number of background events (Sec.~\ref{sec:event_sel}), and the efficiency smeared by resolution in reconstructed bin $i$, respectively. 
$N_{target}$ is the number of argon atoms in the fiducial volume ($2.61 \times 10^{31}$~\cite{Adams:2019iqcINCL}), and $\Phi_{\nu_\mu}$ is the integrated flux of muon neutrinos passing through the detector fiducial volume ($1.17 \times 10^{11}  \text{cm}^{-2}$ for this dataset~\cite{Adams:2019iqcINCL}). 
$(\Delta p_\mu)_i$ is the bin width for bin $i$ in the $p_\mu$ distribution, which was optimised based on a consideration of available statistics and detector resolution.
The other variables ($\cos(\theta_\mu)$, $p_p$, $\cos(\theta_p)$, and $\cos(\theta_{\mu p})$) are handled in a similar way.

To determine if the cross section extraction technique was dependent upon the use of GENIE v2, a sample of simulated events was generated with an alternate GENIE model. The alternate GENIE sample was then treated as observed data, and cross sections extracted using the described technique in the previous paragraphs. The study showed that the cross section extraction is not dependent upon the model used for making the measurement.

Systematic uncertainties (discussed in Sec.~\ref{sec:syst}) on the efficiency and resolution propagate to uncertainties on the final measurement through this procedure, allowing a simple comparison with a smeared theory prediction.




